% ---
% title: Defaults by rating class
% tags: finance, generalized-bayes, total-probability
% ---

\begin{problem}
	\begin{context}
		Banks give every loan an \emph{internal rating}: a grade that summarizes
		how risky the borrower is. For each grade the bank estimates a
		\emph{probability of default} (PD), the chance that the borrower stops
		paying within one year. Better grades have smaller PDs. These numbers
		are used to price loans and to decide how much capital the bank must
		hold.
	\end{context}
	A bank's loan portfolio is split into three rating classes: 50\% of the
	loans are rated A, 30\% are rated B and 20\% are rated C. Within one year a
	loan defaults with probability $0.01$ if it is rated A, $0.03$ if it is
	rated B and $0.08$ if it is rated C. A loan is chosen at random.
	(a) What is the probability that the loan defaults within the year?
	(b) The loan has defaulted. What is the probability that it was rated A,
	B and C, respectively?
	(c) The loan has not defaulted. What is the probability that it was rated
	C?
\end{problem}

\begin{slides}
	\small
	\exstatement
	\pause
	\begin{columns}[T]
		\begin{column}{.47\textwidth}
			\sol Events: $A$, $B$, $C$ = the loan is rated A, B, C;
			$D$ = the loan defaults.
			\pause\\[1mm]
			$A$, $B$, $C$ split $\Om$: every loan is in exactly one class.
			\pause\\[1mm]
			Given: $P(A)=0.5$, $P(B)=0.3$, $P(C)=0.2$;\\
			$P(D\mid A)=0.01$, $P(D\mid B)=0.03$,\\
			$P(D\mid C)=0.08$.
			\uncover<6->{\\\textcolor{teal}{Path: multiply along it},
				e.g.\ $0.5\cdot0.01=0.005$.}
		\end{column}
		\begin{column}{.51\textwidth}
			\pause
			\centering
			\begin{tikzpicture}[
					grow=right, level distance=18mm,
					level 1/.style={sibling distance=12mm},
					level 2/.style={sibling distance=6.2mm},
					every node/.style={font=\scriptsize,inner sep=1pt},
					edge from parent/.style={draw,navy!55},
					lab/.style={fill=white,text=black,inner sep=.8pt,pos=.5}]
				\node {}
				child {node {$C$}
						child {node (c2) {$D^c$} edge from parent node[lab] {$0.92$}}
						child {node[navy] (c1) {$D$} edge from parent node[lab] {$0.08$}}
						edge from parent node[lab] {$0.2$}}
				child {node {$B$}
						child {node (b2) {$D^c$} edge from parent node[lab] {$0.97$}}
						child {node[navy] (b1) {$D$} edge from parent node[lab] {$0.03$}}
						edge from parent node[lab] {$0.3$}}
				child {node {$A$}
						child {node (a2) {$D^c$} edge from parent node[lab] {$0.99$}}
						child {node[navy] (a1) {$D$} edge from parent node[lab] {$0.01$}}
						edge from parent node[lab] {$0.5$}};
				\uncover<6->{\foreach \l/\p in {a1/{0.005},a2/{0.495},b1/{0.009},b2/{0.291},c1/{0.016},c2/{0.184}}
				\node[right=2mm] at (\l) {$\p$};}
			\end{tikzpicture}
		\end{column}
	\end{columns}
\end{slides}

\begin{slides}
	\small
	(a) A path: multiply, e.g.\ $P(DA)=P(D\mid A)P(A)=0.005$. The paths ending
	in $D$ are mutually exclusive: add them.
	\pause
	\[ P(D)=P(D\mid A)P(A)+P(D\mid B)P(B)+P(D\mid C)P(C) \]
	\pause
	\vspace{-3mm}
	\[ P(D)=0.5\cdot0.01+0.3\cdot0.03+0.2\cdot0.08
		=0.005+0.009+0.016=\alert{0.03}. \]
	\pause
	(b) Bayes: one path ending in $D$, divided by all of $P(D)$:
	\[ P(A\mid D)=\frac{P(D\mid A)P(A)}{P(D)}. \]
	\pause
	\vspace{-2mm}
	\[ P(A\mid D)=\frac{0.005}{0.03}=\frac16\approx\alert{0.167},\quad
		\uncover<6->{P(B\mid D)=\frac{0.009}{0.03}=\alert{0.3},}\quad
		\uncover<7->{P(C\mid D)=\frac{0.016}{0.03}=\frac{8}{15}\approx\alert{0.533}.} \]
	\uncover<8->{Check: $0.167+0.3+0.533=1$. A defaulted loan is in one of
		the classes.}
\end{slides}

\begin{slides}
	\small
	(c) Same idea, with the paths ending in $D^c$.
	\pause\\[1mm]
	First: $P(D^c)=1-P(D)=1-0.03=0.97$.
	\pause\\[1mm]
	The path $C\to D^c$: $P(D^c\mid C)P(C)=0.92\cdot0.2=0.184$.
	\pause
	\[ P(C\mid D^c)=\frac{P(D^c\mid C)P(C)}{P(D^c)}=\frac{0.184}{0.97}
		\approx\alert{0.190}. \]
	\pause
	Close to the prior $P(C)=0.2$: almost every loan survives the year, so
	``no default'' tells us little.
	\pause\\[2mm]
	\alert{Lesson:} C is only 20\% of the loans, but it gives more than half
	(53\%) of the defaults. A is half of the loans, but gives only one sixth
	of the defaults.
\end{slides}

\begin{solution}
	\idea The three rating classes split the portfolio into three parts. The
	law of total probability gives the overall default probability in (a). The
	generalized Bayes formula turns ``default given the class'' around into
	``class given default'' in (b) and (c).

	\step{Events and data.} Let $A$, $B$, $C$ be the events that the chosen
	loan is rated A, B, C, and let $D$ be the event that it defaults within the
	year. Every loan is in exactly one class, so $A$, $B$, $C$ are mutually
	exclusive and together they make up $\Om$: they form a partition. The data
	are
	\[ P(A)=0.5,\qquad P(B)=0.3,\qquad P(C)=0.2, \]
	\[ P(D\mid A)=0.01,\qquad P(D\mid B)=0.03,\qquad P(D\mid C)=0.08. \]
	The default probabilities are conditional: $0.08$ is the chance of default
	\emph{for a loan rated C}, not for any loan.

	\step{The tree.} The first level is the rating class, the second level is
	default or not. The probability at the end of a path is the product of the
	probabilities along it. This is the multiplication rule; for example, the
	path through $A$ to $D$ gives
	\[ P(DA)=P(D\mid A)P(A)=0.01\cdot0.5=0.005. \]
	The non-default probabilities are the complements:
	\[ P(D^c\mid A)=0.99,\qquad P(D^c\mid B)=0.97,\qquad P(D^c\mid C)=0.92. \]
	\begin{center}
		\begin{tikzpicture}[
				grow=right, level distance=24mm,
				level 1/.style={sibling distance=17mm},
				level 2/.style={sibling distance=8mm},
				every node/.style={font=\small},
				edge from parent/.style={draw,gray}]
			\node {$\Om$}
			child {node {$C$}
					child {node (c2) {$D^c$} edge from parent node[below,font=\scriptsize] {$0.92$}}
					child {node[navy] (c1) {$D$} edge from parent node[above,font=\scriptsize] {$0.08$}}
					edge from parent node[below left,font=\scriptsize] {$0.2$}}
			child {node {$B$}
					child {node (b2) {$D^c$} edge from parent node[below,font=\scriptsize] {$0.97$}}
					child {node[navy] (b1) {$D$} edge from parent node[above,font=\scriptsize] {$0.03$}}
					edge from parent node[above,font=\scriptsize,pos=.6] {$0.3$}}
			child {node {$A$}
					child {node (a2) {$D^c$} edge from parent node[below,font=\scriptsize] {$0.99$}}
					child {node[navy] (a1) {$D$} edge from parent node[above,font=\scriptsize] {$0.01$}}
					edge from parent node[above left,font=\scriptsize] {$0.5$}};
			\foreach \l/\p in {a1/{0.5\cdot0.01=0.005},a2/{0.5\cdot0.99=0.495},
					b1/{0.3\cdot0.03=0.009},b2/{0.3\cdot0.97=0.291},
					c1/{0.2\cdot0.08=0.016},c2/{0.2\cdot0.92=0.184}}
			\node[right=2mm,font=\scriptsize] at (\l) {$\p$};
		\end{tikzpicture}
	\end{center}
	As a check, the six path probabilities add up to $1$.

	\step{(a)} A defaulted loan belongs to exactly one class, so $D$ splits
	into the three mutually exclusive pieces $DA$, $DB$, $DC$: the three paths
	that end in $D$. Their probabilities add up (law of total probability):
	\[ \begin{aligned}
			P(D) & =P(D\mid A)P(A)+P(D\mid B)P(B)+P(D\mid C)P(C) \\
			     & =0.01\cdot0.5+0.03\cdot0.3+0.08\cdot0.2
			=0.005+0.009+0.016=0.03.
		\end{aligned} \]
	So 3\% of the loans default. This is a weighted average of the three
	default probabilities, with the class shares as weights.

	\step{(b)} We need $P(A\mid D)$, $P(B\mid D)$ and $P(C\mid D)$. By the
	definition of conditional probability,
	\[ P(A\mid D)=\frac{P(AD)}{P(D)}. \]
	The numerator is one path of the tree, the denominator is the sum of all three
	default paths from (a). This is the generalized Bayes formula: for a
	partition $F_1,\dots,F_n$,
	\[ P(F_j\mid D)=\frac{P(D\mid F_j)P(F_j)}{\sum_{i=1}^{n}P(D\mid F_i)P(F_i)}. \]
	Here the partition is $A$, $B$, $C$:
	\[ P(A\mid D)=\frac{0.005}{0.03}=\frac16\approx0.167,\qquad
		P(B\mid D)=\frac{0.009}{0.03}=0.3,\qquad
		P(C\mid D)=\frac{0.016}{0.03}=\frac{8}{15}\approx0.533. \]
	The three numbers add up to $1$ (exactly: $\frac{5}{30}+\frac{9}{30}+
		\frac{16}{30}=1$), as they should: a defaulted loan was rated A, B or C.

	\step{(c)} Now the condition is $D^c$, the loan did not default. Its
	probability is the complement of (a): $P(D^c)=1-0.03=0.97$ (or add the
	three paths ending in $D^c$: $0.495+0.291+0.184=0.97$). The path through
	$C$ to $D^c$ has probability $P(D^c\mid C)P(C)=0.92\cdot0.2=0.184$. So
	\[ P(C\mid D^c)=\frac{P(D^c\mid C)P(C)}{P(D^c)}=\frac{0.184}{0.97}
		\approx0.190. \]
	\begin{answer}
		(a) $P(D)=0.03$ \quad (b) $P(A\mid D)=\frac16\approx0.167$,
		$P(B\mid D)=0.3$, $P(C\mid D)=\frac{8}{15}\approx0.533$ \quad
		(c) $P(C\mid D^c)=\frac{0.184}{0.97}\approx0.190$
	\end{answer}
	\begin{remark}
		Class C makes up only 20\% of the loans, but more than half (53\%) of the
		defaults come from it. Class A is half of the portfolio, yet it causes
		only one sixth of the defaults. Learning that a loan has \emph{not}
		defaulted changes very little: $P(C\mid D^c)\approx0.190$ is close to
		the prior $0.2$, because almost every loan survives the year.
		\textbf{Common mistake:} answering $0.08$ for $P(C\mid D)$. That number
		is $P(D\mid C)$, the default probability of a C loan; (b) asks the other
		way around.
	\end{remark}
\end{solution}
